% 原书物理页=印刷页60--63；2.4 全节.（2.4.3 段 Weyl 定理证明的后半部分与其后的推论实际印于原书第 63 页上部、2.5 节标题之前，属本节内容，一并译出.）
\section{半单 Lie 代数}\label{sec:lie-2-4}

设 $\mathfrak h,\mathfrak h'$ 是 $K$-Lie 代数 $\mathfrak g$ 的两个可解理想，则它们的和 $\mathfrak h+\mathfrak h'$ 仍是可解的，因为它是直和 $\mathfrak h\oplus\mathfrak h'$ 的一个商代数. 这使得我们能够定义 $\mathfrak g$ 的\newterm{lie-radical}{根基}{radical} $\operatorname{rad}(\mathfrak g)$，即 $\mathfrak g$ 的\emph{最大可解理想}.
\begin{definition}\label{def:lie-2-4-1}
称一个 $K$-Lie 代数 $\mathfrak g$ 是
\begin{itemize}[label=--]
\item \newterm{lie-simple-lie-algebra}{单}{simple}的，如果 $\dim(\mathfrak g)>1$ 且它仅有的理想是 $\{0\}$ 与 $\mathfrak g$；
\item \newterm{lie-semisimple}{半单}{semi-simple}的，如果 $\operatorname{rad}(\mathfrak g)=0$（即 $\mathfrak g$ 没有非零的阿贝尔理想）；
\item \newterm{lie-reductive}{约化}{reductive}的，如果 $\operatorname{rad}(\mathfrak g)=\mathfrak z(\mathfrak g)$.
\end{itemize}
\end{definition}
\begin{remark*}
``单''的定义中条件 $\dim(\mathfrak g)>1$ 是为了排除``平凡''代数 $\mathfrak g=K$，后者仅有的理想是 $\{0\}$ 与 $\mathfrak g$. 在这些定义下，我们有
\[
\mathfrak g\ \text{单}\Rightarrow\mathfrak g\ \text{半单}\Rightarrow\mathfrak g\ \text{约化}
\]
\end{remark*}
\begin{exercise*}
\begin{enumerate}[label=\roman*)]
\item 证明 $\mathfrak g$ 单 $\Rightarrow\mathfrak g$ 半单.
\item 证明对任意 $\mathfrak g$，$\mathfrak g/\operatorname{rad}(\mathfrak g)$ 是半单的.
\item 证明半单 Lie 代数的直和是半单的.
\end{enumerate}
\end{exercise*}
我们将会看到，所有经典 Lie 代数都是约化的，而且一旦其中心平凡，它们甚至是单的.
\begin{numpara}[用 Killing 型刻画]\label{np:lie-2-4-2}
从现在起我们假定 $K$ 代数闭且特征为零.
\begin{lemma*}
设 $(V,r)$ 是 $\mathfrak g$ 的一个不可约表示. 则 $\operatorname{rad}(\mathfrak g)$ 中任何 $X$ 都以一个纯量作用，从而 $r([\mathfrak g,\operatorname{rad}(\mathfrak g)])=0$.
\end{lemma*}
\begin{proof}
由 Lie 定理，$V$ 中存在 $\operatorname{rad}(\mathfrak g)$ 的一个特征向量. 把与之相伴的``特征值''记作 $\lambda:\operatorname{rad}(\mathfrak g)\longrightarrow K$，并记 $V_{\lambda}$ 为相应的特征空间. 我们只需证明 $V_{\lambda}$ 在 $\mathfrak g$ 下稳定. 事实上，对一切 $Y\in\mathfrak g$、$X\in\operatorname{rad}(\mathfrak g)$ 与 $v\in V_{\lambda}$，有
\[
XYv=[X,Y]v+YXv=\lambda([X,Y])v+\lambda(X)Yv.
\]
与 Lie 定理的证明中一样，可以证明 $\lambda([X,Y])=0$，于是得证.
\end{proof}
\begin{corollary*}
若 $\mathfrak g$ 容许一个表示 $(V,r)$ 使得双线性型 $B_{r}$ 非退化，则 $\mathfrak g$ 是约化的.
\end{corollary*}
\begin{proof}
由上面的引理，理想 $[\mathfrak g,\operatorname{rad}(\mathfrak g)]$ 在 $\mathfrak g$ 的任何不可约表示上平凡作用. 于是取 $V$ 的一个由子表示组成的递增序列 $0\subset V_{1}\subset V_{2}\subset\cdots\subset V_{l}=V$，其相邻商 $V_{i}/V_{i-1}$ 均不可约. 因而对一切 $Y\in\mathfrak g$ 有 $r(Y)(V_{i})\subset V_{i}$，而对一切 $X\in[\mathfrak g,\operatorname{rad}(\mathfrak g)]$ 有 $r(X)(V_{i})\subset V_{i-1}$. 因此 $r(X)r(Y)$ 是幂零的，从而迹为零，于是 $[\mathfrak g,\operatorname{rad}(\mathfrak g)]$ 含于 $B_{r}$ 的核中. 由于 $B_{r}$ 被假定非退化，故 $[\mathfrak g,\operatorname{rad}(\mathfrak g)]=0$，这等价于 $\operatorname{rad}(\mathfrak g)\subset\mathfrak z(\mathfrak g)$. 反向包含是重言式的.
\end{proof}
\begin{example*}
把这个推论应用于一个经典 Lie 代数的标准表示，我们发现 $\mathfrak{gl}_n(K)$ 与 $\mathfrak u(n)$（当 $K=\mathbb R$ 时）是约化的，而 $\mathfrak{sl}_n(K)$、$\mathfrak{so}_n(K)$（当 $n>2$ 时）、$\mathfrak{sp}_{2n}(K)$、$\mathfrak{su}(n)$（当 $K=\mathbb R$ 时）是约化的，甚至是半单的，因为（习题）它们的中心为零.
\end{example*}
\begin{theorem*}[Cartan]
一个 $K$-Lie 代数 $\mathfrak g$ 是半单的当且仅当 Killing 型 $B_{\mathrm{ad}}$ 非退化.
\end{theorem*}
\begin{proof}
设 $B_{\mathrm{ad}}$ 非退化. 这显然蕴含 $\operatorname{Ker}(\operatorname{ad})=\mathfrak z(\mathfrak g)=0$，并且由上面的推论，它也蕴含 $\mathfrak g$ 是约化的. 故 $\mathfrak g$ 是半单的.
反之，设 $\mathfrak g$ 半单. 那么 $\mathfrak z(\mathfrak g)=0$，故 $\operatorname{ad}:\mathfrak g\longrightarrow\mathfrak{gl}(\mathfrak g)$ 是单射. 设 $\mathfrak h=\mathfrak g^{\perp}$ 为 $B_{\mathrm{ad}}$ 的核. 由 Cartan 判据，$\mathfrak h$ 是可解的，故 $\mathfrak h=0$.
\end{proof}
\begin{theorem*}
一个半单 $K$-Lie 代数 $\mathfrak g$ 是单理想的直和. 特别地，$\mathfrak g=[\mathfrak g,\mathfrak g]$.
\end{theorem*}
\begin{proof}
设 $\mathfrak h$ 是 $\mathfrak g$ 的一个理想，$\mathfrak h^{\perp}$ 是它关于 $\mathfrak g$ 的 Killing 型的正交补. 那么 $\mathfrak h\cap\mathfrak h^{\perp}$ 的 Killing 型（它是 $\mathfrak g$ 的 Killing 型的限制）为零，故 $\mathfrak h\cap\mathfrak h^{\perp}$ 是一个可解理想，从而为零. 由于 $B_{\mathrm{ad}}$ 非退化，由此推出 $\mathfrak g=\mathfrak h\oplus\mathfrak h^{\perp}$，且 $\mathfrak h$ 与 $\mathfrak h^{\perp}$ 的 Killing 型都非退化. 于是对维数作归纳即可得结论.
\end{proof}
\end{numpara}
\begin{numpara}[用表示刻画]\label{np:lie-2-4-3}
仍设 $K$ 代数闭且特征为零.
\begin{theorem*}[Weyl]
一个 $K$-Lie 代数 $\mathfrak g$ 是半单的当且仅当它的所有（有限维）表示都完全可约.
\end{theorem*}
Weyl 对该定理的原始证明（当 $K=\mathbb C$ 时）用到了他的\newterm{lie-unitary-trick}{幺正技巧}{unitary trick}以及一个作为紧 Lie 群的 Lie 代数出现的实形式的存在性，正如我们对 $\mathfrak{sl}_n(\mathbb C)$ 用实形式 $\mathfrak{su}(n)$ 与紧群 $\mathrm{SU}(n)$ 所做的那样. 下面给出一个归功于 Casimir 的更直接的证明.
\begin{proof}
若诸表示都完全可约，则特别地伴随表示也如此. 写 $\mathfrak g=\bigoplus_{i}\mathfrak g_{i}$，其中诸 $\mathfrak g_{i}$ 是 $(\mathfrak g,\operatorname{ad})$ 的不可约子表示. 那么诸 $\mathfrak g_{i}$ 都是理想（因为它们在 $\operatorname{ad}_{\mathfrak g}$ 下稳定），并且是单的或维数为 $1$（因为它们不可约）. 然而 $\mathfrak g$ 不能有维数为 $1$ 的直因子，因为我们已经看到一维 Lie 代数 $K$ 的表示未必完全可约. 故 $\mathfrak g$ 是半单的.
现在设 $\mathfrak g$ 半单. 设 $(V,r)\in\operatorname{Rep}_{K}(\mathfrak g)$，$W\subset V$ 是一个子表示；我们想找到 $W$ 的一个在 $\mathfrak g$ 下稳定的补. 为此，只需找到一个 $\varphi\in\operatorname{Hom}_{K\mathfrak g}(V,W)$ 使得 $\varphi|_{W}=\lambda\,\mathrm{id}_{W}$，其中 $\lambda\neq 0$. 因为事实上我们将有 $V=W\oplus\operatorname{Ker}(\varphi)$，且 $\operatorname{Ker}(\varphi)$ 稳定.
\emph{第一步：化归到 $V/W$ 为单位表示的情形.}\ 考虑内 Hom 表示 $H=\mathcal{H}om_{K}(V,W)$ 及其子空间 $A:=\{\varphi\in H,\ \varphi|_{W}\in K\,\mathrm{id}_{W}\}$ 与 $B:=\{\varphi\in A,\ \varphi|_{W}=0\}$. 若 $\varphi\in A$，则可看到对一切 $X\in\mathfrak g$ 有 $(X\varphi)(w)=X\varphi(w)-\varphi(Xw)=X\lambda w-\lambda Xw=0$，故 $X\varphi\in B$. 特别地，$A$ 与 $B$ 都是 $H$ 的子表示，且 $B/A$ 是单位表示. 假设 $B$ 容许一个补，也就是说存在一条直线 $D\subset A^{\mathfrak g}$ 使得 $B\oplus D=A$. 那么，如上所释，我们有 $W\oplus\operatorname{Ker}(\varphi)=V$.
\emph{第二步：Casimir 算子与归纳.}\ 现在我们假设 $V/W$ 是单位表示，并对 $\dim(V)$ 作归纳来证明 $W$ 存在一条稳定补直线. 取定一个由子表示组成的递增序列 $0=V_{0}\subset V_{1}\subset\cdots\subset V_{l-1}=W\subset V_{l}=V$，使得 $V_{i}/V_{i-1}$ 对一切 $i=1,\cdots,l$ 不可约.
若所有 $V_{i}/V_{i-1}$ 都同构于单位表示，则我们面对的是 $V$ 的一个完全旗 $\mathcal V$ 使得 $r(\mathfrak g)\subset\mathfrak b^{\mathrm{nilp}}(\mathcal V)$. 但由于对一切 $n$ 有 $\mathfrak g=[\mathfrak g,\mathfrak g]=C^{n}\mathfrak g$，取 $n=\dim(V)$ 便推出 $r(\mathfrak g)=r(C^{n}\mathfrak g)=C^{n}\mathfrak b^{\mathrm{nilp}}(\mathcal V)=0$，故 $V$ 是平凡表示，从而 $W$ 的任何补直线都是稳定的.
否则，存在 $i$ 使得 $(U,r_{U}):=V_{i}/V_{i-1}$ 异于单位表示. 此时像 $r_{U}(\mathfrak g)$ 非零，于是可以选取 $\mathfrak g$ 的一个单直因子 $\mathfrak g_{U}$ 使得 $r_{U}(\mathfrak g_{U})\neq 0$. 于是 $\mathfrak g_{U}\xrightarrow{\ \sim\ }r_{U}(\mathfrak g_{U})$，故 $r_{U}(\mathfrak g_{U})$ 是单的，而与 $r_{U}$ 相配的双线性型 $B_{r_{U}}$ 非零（否则由 Cartan 判据将推出 $r_{U}(\mathfrak g_{U})$ 可解）. 它的核是 $\mathfrak g_{U}$ 的一个真理想，故为零，最终 $B_{r_{U}}$ 在 $\mathfrak g_{U}$ 上非退化. 设 $C_{U}$ 为与之相伴的 Casimir 算子，它属于 $\mathcal U\mathfrak g_{U}$ 的中心，从而属于 $\mathcal U\mathfrak g$ 的中心（事实上 $\mathcal U(\mathfrak h\oplus\mathfrak h')=\mathcal U\mathfrak h\otimes_{K}\mathcal U\mathfrak h'$）. 那么 $r(C_{U})$ 是 $V$ 的一个与 $\mathfrak g$ 的作用交换的自同态，故 $\mathfrak g$ 稳定它的各广义特征子空间. 我们有 $r(C_{U})(V)\subset W$，故 $0$ 是一个特征值，且广义特征空间 $V_{0}:=\operatorname{Ker}(r(C_{U})^{n})$ 不含于 $W$. 设 $V^{0}$ 为 $r(C_{U})$ 的其余广义特征子空间之和. 于是 $V=V^{0}\oplus V_{0}$ 是由稳定子空间组成的直和. 若能证明 $V^{0}\neq 0$，则可把归纳假设应用于 $V_{0}$. 因此我们必须证明 $r(C_{U})$ 具有一个非零特征值. 为此只需验证 $r_{U}(C_{U})$ 具有一个非零特征值. 注意，由 Schur 引理，$r_{U}(C_{U})$ 是一个位似. 因此只需证明它的迹非零. 然而 $\operatorname{tr}(r_{U}(C_{U}))=\sum_{i=1}^{\dim\mathfrak g_{U}}\operatorname{tr}(r_{U}(X_{i})r_{U}(X_{i}^{*}))$，其中 $X_{i}$ 是 $\mathfrak g_{U}$ 的一组基，而 $X_{i}^{*}$ 是关于双线性型 $B_{r_{U}}(X,Y)=\operatorname{tr}(r_{U}(X)r_{U}(Y))$ 的对偶基. 于是 $\operatorname{tr}(r_{U}(C_{U}))=\dim\mathfrak g_{U}\neq 0$.
\end{proof}
\begin{corollary*}
一个 $K$-Lie 代数 $\mathfrak g$ 是约化的当且仅当 $\mathfrak g=\mathfrak z(\mathfrak g)\oplus[\mathfrak g,\mathfrak g]$，其中 $[\mathfrak g,\mathfrak g]$ 半单.
\end{corollary*}
\begin{proof}
该条件的充分性是显然的. 我们来证明它是必要的. 设 $\mathfrak g$ 约化. 那么伴随表示 $\operatorname{ad}:\mathfrak g\longrightarrow\mathfrak{gl}(\mathfrak g)$ 可通过 $\operatorname{ad}(\mathfrak g)=\mathfrak g/\mathfrak z(\mathfrak g)$ 分解，而后者是半单的. 因而它是完全可约的. 由于 $\mathfrak z(\mathfrak g)\subset\mathfrak g$ 在 $\operatorname{ad}(\mathfrak g)$ 下稳定，故存在分解 $\mathfrak g=\mathfrak z(\mathfrak g)\oplus\mathfrak g_{1}\oplus\cdots\oplus\mathfrak g_{k}$，其中诸 $\mathfrak g_{i}$ 不可约，从而是单的. 由此推出 $[\mathfrak g,\mathfrak g]=[\mathfrak g_{1},\mathfrak g_{1}]\oplus\cdots\oplus[\mathfrak g_{k},\mathfrak g_{k}]=\mathfrak g_{1}\oplus\cdots\oplus\mathfrak g_{k}$.
\end{proof}
\end{numpara}
